\section{Supplied-Record Audit}
\label{app:record_audit}

The supplied feature-level dataset describes 40 de-identified healthy participants (20
female and 20 male at birth), with ten participants in each of four age bands
(18--29, 30--44, 45--59, and 60--74 years). Each participant contributed five
visits, yielding 200 sessions and 800 B/E/F/A component evaluations. Visits were
separate acquisitions spaced exactly seven days apart (160 successive-visit
intervals; median and range, 7 days). The
supplied participant metadata mark inclusion as passed, consent as obtained,
and ethics review as approved. Committee identifier, approval documentation,
and an independently auditable acquisition log are not included; these fields
are source assertions rather than verified approvals. The supplied session metadata describe collection in a controlled
indoor simulated-home environment using a Raspberry Pi Camera Module 3 at
1280$\times$720 and a
target rate of 15 frames/s. We treated a module observation as analyzable when its
state was neither \textit{insufficient} nor \textit{skipped}. For each
selected module output, Table~\ref{tab:human_repeatability} reproduces a
one-way single-measure ICC(1,1), computed from the analyzable observations, a
participant-bootstrap 95\% interval, and the residual within-participant SD
normalized numerically by the corresponding engineering reference. For F, the
unresolved source-unit mismatch prevents interpreting that ratio as variation
relative to the implemented threshold. Because four
observations per module failed quality gating, the ICC calculation used 196
observations from 40 participants, with four or five repeats per participant.
As a sensitivity analysis, ICC(2,1) used the 36 participants with all five
visits for the relevant module.

\begin{table}[htbp]
\centering
\caption{Descriptive reanalysis of supplied records from 40 participants; baseline provenance and F units remain unresolved.}
\label{tab:human_repeatability}
\footnotesize
\setlength{\tabcolsep}{3pt}
\renewcommand{\arraystretch}{1.03}
\begin{tabular}{@{}lrrrr@{}}
\toprule
\textbf{Module} & \textbf{ICC(1,1)} & \textbf{95\% CI} & \textbf{ICC(2,1)} & \textbf{SD/ref.} \\
\midrule
B & 0.1404 & 0.013--0.498 & 0.1440 & 8.25\% \\
E & 0.2662 & 0.089--0.817 & 0.2573 & 9.50\% \\
F & 0.7837 & 0.649--0.883 & 0.7882 & 9.81\% \\
A & 0.6788 & 0.429--0.880 & 0.6579 & 9.30\% \\
\bottomrule
\end{tabular}
\end{table}

For every module, 195/200 evaluations were negative, 1/200 was positive, and
4/200 were insufficient, giving an analyzable rate of 98\%. Mean quality scores
were 0.933 for B, 0.928 for E, 0.931 for F, and 0.931 for A. The positive state
in each module coincided with a raw metric crossing its engineering reference;
the four insufficient states coincided with the module-specific quality-failure
reason. These supplied healthy-control records describe selected-output variation and
recorded quality yield, but do not estimate clinical
sensitivity or specificity. The B and E numeric ICC estimates were lower than F and A. These
relative-variation statistics do not resolve measurement provenance or establish
absolute repeatability of the implemented modules.
The E status counts were recorded as an offline
engineering-reference audit. Unvalidated E thresholds remain disabled in the
released runtime.

The supplied age--sex cross-tabulation has ten women in each of the
18--29 and 45--59 bands and ten men in each of the other two bands; age-band
and sex effects cannot be separated. More fundamentally, the exported B change
scores have no linked five-window enrollment records, so overlap with evaluation
visits and baseline freezing cannot be audited. F is labeled
\textit{shoulder-width ratio} in the source whereas the method uses interocular
normalization. We preserve the source values and do not infer a conversion.
The reported F SD/reference is therefore only a source-numeric calculation.
Only one selected output per module is exported; the full multi-feature runtime
decisions and quality gates cannot be reconstructed from these tables. These
results are descriptive reanalysis pending provenance reconciliation, not a
validated repeatability study of the released implementation.

\begin{figure}[htbp]
  \centering
  \includegraphics[width=0.88\columnwidth]{repeatability-icc.pdf}
  \caption{Selected-output repeatability. Intervals resample participants;
  quality-gate failures remain in yield summaries but are excluded from ICC.}
  \label{fig:repeatability_icc}
\end{figure}

ICC is sensitive to the restricted between-participant range of a healthy cohort
and therefore does not by itself establish agreement. The within-participant
SD/reference values are engineering-normalized dispersion, not standard clinical
repeatability coefficients. Future reporting should add Bland--Altman limits,
standard error of measurement, minimal detectable change, and sensitivity
analyses that retain quality failures. The one engineering-rule positive per 200
evaluations is a threshold-crossing proportion, not a clinically adjudicated
false-positive rate.

